\section{A.3. Cutting and welding with the necessary regularity}
\begin{lemma}[Proved: the boundary-chart exponent in 67/055--056]
\label{A3-boundary-chart}
Let $f$ be a nonconstant holomorphic germ at zero with $f(0)=0$, real
values on the real axis, and $\operatorname{Im} f(z)\geq0$ for $z$ in
the upper half-disc. Then $f'(0)>0$. In particular the germ is locally
biholomorphic, as required by Th\'eor\`eme 1 of cote 67/055--056.
\end{lemma}
\begin{proof}
Write $f(z)=az^m+O(z^{m+1})$, with $m\geq1$ and nonzero real $a$;
the coefficients are real by Schwarz reflection. For any fixed
$0<\theta<\pi$, dividing $\operatorname{Im}f(re^{i\theta})\geq0$
by $r^m$ and letting $r\downarrow0$ gives $a\sin(m\theta)\geq0$.
If $m\geq2$, this sine takes both positive and negative values on
$(0,\pi)$, a contradiction. Thus $m=1$, and the same inequality at
$\theta=\pi/2$ gives $a>0$. The holomorphic inverse theorem applies.
\end{proof}

\begin{lemma}[Proved: the continuous-function removability of Lemma 3(c)]
\label{A3-rectifiable}
Let $U\subset\C$ be open and let $C$ be a relatively closed
rectifiable curve in $U$, or locally a finite union of rectifiable
arcs. Every continuous complex-valued function on $U$ which is
holomorphic on $U\setminus C$ is holomorphic throughout $U$.
In particular the continuous-function assertion of 67/061 is valid;
this does not by itself assert existence of a welding for an
arbitrary interval homeomorphism.
\end{lemma}
\begin{proof}
It is enough to work in a rectangle $R$ compactly contained in a
coordinate neighborhood where the relevant curves have total
finite length $L$. Subdivide $R$ into rectangles with both side
lengths between $\delta/2$ and $\delta$. Divide each curve into
arclength pieces of length at most $\delta$. Each piece lies in
a disc of radius $\delta$, meeting a uniformly bounded number of
the grid cells. Hence at most $C_0(1+L/\delta)$ cells meet the
curves, with $C_0$ independent of $\delta$ (the finite number of
arcs is absorbed in the constant).

Write $\omega_f$ for the modulus of continuity of $f$ on $R$.
The integral of $f\,dz$ around a cell avoiding $C$ vanishes by
Cauchy's theorem. On any other cell $Q$, subtract the constant
$f(z_Q)$ for a chosen $z_Q\in Q$; its boundary integral is zero,
and therefore
\[
 \left|\int_{\partial Q}f\,dz\right|
       \leq4\delta\,\omega_f(\sqrt2\delta).
\]
Sum over all cells. Interior edges cancel with opposite
orientations, so
\[
 \left|\int_{\partial R}f\,dz\right|
       \leq4C_0(L+\delta)\omega_f(\sqrt2\delta)
       \longrightarrow0.
\]
Thus every such rectangle has zero boundary integral. The
rectangle form of Morera's theorem, proved by defining local
primitives along horizontal and vertical segments, shows that
$f$ is holomorphic. The argument is local and proves the stated
finite-union version as well.
\end{proof}

\begin{theorem}[Known with proof of the form used here: circle welding]\label{A3-welding}
Let $X_1,X_2$ be bordered Riemann surfaces with analytic boundary, and identify selected boundary circles by orientation-reversing homeomorphisms that are quasisymmetric in analytic boundary coordinates. For finitely many such identifications, the topological quotient has a conformal structure for which the inclusions of the interiors are holomorphic and the seams are locally quasiconformal images of line segments. It is unique up to the unique conformal isomorphism restricting to the identity on every original piece. More generally, whenever a welding exists, the same uniqueness conclusion holds if its seam is conformally removable.
\end{theorem}
\begin{proof}
Beurling--Ahlfors \cite[Theorem 1, p.~126]{ABeurlingAhlfors1956} gives a quasiconformal extension of each local boundary correspondence to a half-disc; changing to circle coordinates uses a conformal map between a half-plane and a disc. On a collar of a seam, extend the boundary identification into one side and keep the other side fixed. This supplies a quasiconformal topological coordinate on the quotient, with a measurable Beltrami coefficient of norm strictly less than one. There are only finitely many circles, so the dilatations have a common finite bound. Integrating this Beltrami coefficient by the measurable Riemann mapping theorem supplies a holomorphic coordinate on each quotient neighborhood. On overlaps the comparison has zero Beltrami coefficient on both open sides and is quasiconformal across the straightened seam; hence it is conformal almost everywhere and therefore holomorphic. These charts give the claimed conformal structure and the seam is a quasiarc locally.

For completeness, a quasiarc is locally conformally removable. Straighten it with a quasiconformal map. A homeomorphism conformal off the original arc becomes quasiconformal off a line segment with a uniform dilatation bound. On almost every parallel line it is absolutely continuous on the two complementary intervals, and continuity glues those intervals; local area bounds give the required square-integrable weak derivatives. It is consequently quasiconformal across the line. Returning to the original coordinate makes it quasiconformal everywhere, with zero Beltrami coefficient off a set of area zero. Quasiconformal maps preserve area-zero sets, so the coefficient vanishes almost everywhere, proving conformality. This argument also works after covering a seam circle by finitely many neighborhoods.

If two welded structures are given, their topological comparison is the identity on the pieces and is conformal off the seam. Removability gives conformality across it. Any comparison fixing the open pieces is already unique as a continuous map. This proves both the stated uniqueness and its general removable-seam version. Bishop's introductory discussion and quasisymmetric-welding argument \cite[Sections 1 and 4, pp.~613--615 and 630--632]{ABishop2007} give another construction of the circle case.
\end{proof}

\begin{proposition}[Proved under explicit vertex compatibility: graph welding]\label{A3-graph}
Theorem \ref{A3-welding} extends to finite systems of orientation-reversing boundary-arc identifications, provided that the topological quotient is an oriented surface and the identifications are quasiconformal-compatible at the vertices: in a neighborhood of each vertex there is an orientation-preserving homeomorphism to a standard finite star of sectors which is quasiconformal on each incident open sector with one common finite bound. Cutting and welding are then inverse equivalences between these marked seam surfaces and their cut surfaces with pairing data. Smooth arc identifications with compatible nonzero derivatives and finite nonzero sector angles satisfy this condition.
\end{proposition}
\begin{proof}
Away from the finitely many vertices apply Theorem \ref{A3-welding}. At a vertex use the sector model in the hypothesis and transfer the measurable complex structure from every open sector. Its Beltrami coefficient has one bound strictly below one; assign any value on the finitely many rays and at the vertex, which have area zero. Integrate this coefficient on the full disc, or reflect it and integrate on the double of a half-disc at a boundary vertex. This constructs the vertex chart directly. The seam is locally a finite union of quasiarcs. Their open edges are removable by the previous proof, and their finitely many endpoints are removable afterwards. Cutting recovers the given prime-end sides and welding recovers the quotient, including all incidence and pairing data.

Here is an explicit verification of the smooth sufficient condition. Straighten the finitely many smooth corners by local bi-Lipschitz coordinates, with sector angles $\alpha_j>0$. Put $A=\sum_j\alpha_j$ and $c=\Omega/A$, where $\Omega=2\pi$ at an interior vertex and $\Omega=\pi$ at a boundary vertex. Use target angles $c\alpha_j$ and choose one smooth radial parameter on each quotient ray. Its two pullbacks to the boundary rays of sector $j$ are functions $s_{j,-}(r),s_{j,+}(r)$ vanishing at zero, with derivatives bounded above and away from zero. On that sector, with $u=\theta/\alpha_j$, define
\[
 R_j=\exp\!\left(c\bigl((1-u)\log s_{j,-}(r)
                   +u\log s_{j,+}(r)\bigr)\right),
 \qquad \Theta_j=\Theta_{j,0}+c\theta .
\]
The maps agree along each paired ray, where their radial coordinate is the same quotient parameter raised to the \emph{common} power $c$. In logarithmic polar coordinates,
\[
 \partial_{\log r}\log R_j
 =c\left((1-u)\frac{r s'_{j,-}}{s_{j,-}}
                 +u\frac{r s'_{j,+}}{s_{j,+}}\right)>0,
 \qquad
 \partial_\theta\log R_j
 =\frac c{\alpha_j}\log\frac{s_{j,+}}{s_{j,-}}.
\]
The first quantity is bounded above and away from zero, and the second is bounded, because $s_{j,\pm}(r)/r$ tend to positive constants. The angular derivative is $c$ and its radial derivative is zero. Thus these matrices and their inverses have uniform finite bounds, proving quasiconformality. There are finitely many sectors, so one bound works on the whole vertex star. This construction also explains why unrelated power normalizations on the different sectors would not by themselves prove compatibility.
\end{proof}

\begin{theorem}[Known with derivation: slit uniformization and its corrected converse]\label{A3-slit}
Let $V$ be a Jordan domain and let $I$ be a Jordan arc joining $a\in V$ to $b\in\partial V$, with $I\setminus\{b\}\subset V$. No rectifiability is required for the conformal uniformization of the cut domain $V\setminus I$ by a disc. Its prime-end boundary consists of the outer boundary and two copies of the slit, with a single tip and two copies of its root. Specifying three positively cyclically ordered prime ends fixes the uniformizing map uniquely.

For the converse in Lemma 4B and the Scholie of 67/072--073, take the three source points $1,\zeta,\bar\zeta$, $\zeta=e^{2\pi i/3}$, and the arcs $\beta=[1,\zeta]$, $\bar\beta=[1,\bar\zeta]$ in the source circle. Let $h:\beta\to\bar\beta$ identify the tip $1$ with itself and the two roots with each other. Assume the induced order-reversing involution on the combined arc $\beta\cup\bar\beta$ is quasisymmetric, with the endpoint compatibility of Proposition \ref{A3-graph}. Then there is a unique continuous map
\[
 f:\overline\D\to\overline\D
\]
which is conformal and injective on $\D$, has image $\D\setminus I$ there for a quasislit $I$ from $0$ to $1$, sends the source tip $1$ to $0$ and the source roots $\zeta,\bar\zeta$ to $1$, and realizes $h$ as its bank identification. Conversely, a quasiconformal image of the standard slit pair gives such pairing data. For an arbitrary slit whose welding exists, uniqueness holds whenever the slit is removable; one must not assert that every quasisymmetric welding produces a rectifiable slit.
\end{theorem}
\begin{proof}
The cut prime-end closure is a closed topological disc: this can be constructed by taking two banks of the arc and separating their root endpoints. A Jordan arc is locally connected, as is the union with the Jordan boundary. The Riemann mapping theorem and Carath\'eodory's prime-end boundary extension therefore give the claimed continuous uniformization and boundary description; the precise completion and locally connected boundary statements, with proofs, are \cite[Theorems 15.9, 16.6 and 16.7, pp.~15-5 and 16-2--16-4]{AMilnor1991}. The automorphism group of $\D$ acts simply transitively on positively ordered triples of distinct boundary points; this proves the first uniqueness statement.

For the converse form the topological quotient by the bank pairing. It is a disc with a marked interior tip and a marked boundary root. The hypotheses allow Proposition \ref{A3-graph} to furnish its conformal structure with a removable quasislit. Uniformize this quotient by $\D$, send the interior tip to zero, and use the remaining rotation to send the root to one. The composite from the original disc has all the required properties. A conformal self-map of $\D$ fixing zero and one is the identity, proving uniqueness. Conversely, lift a quasiconformal equivalence of slit pairs to their prime-end cut discs and conjugate by the normalized Riemann maps. The resulting quasiconformal disc map has a quasisymmetric boundary map by Beurling--Ahlfors. Conjugating the standard bank pairing gives the claimed involution and its vertex compatibility.

The distinction at the tip is necessary for a clean quasislit statement: an interval map can be quasisymmetric on each individual bank while its doubled involution fails quasisymmetry at their common tip. Also, quasisymmetric weldings can have nonrectifiable seams, so the word ``rectifiable'' cannot be retained as the converse's conclusion. For the doubtful transcription $f(\omega\mid0)=0$ on scan 073, the reading $f(0)=0$ is incompatible with the asserted image $\D\setminus I$ and $0\in I$; the correctly typed normalization is the boundary tip mapping to zero.
\end{proof}

\begin{corollary}[Proved under the actual geodesic-cut hypotheses: L5]\label{A3-L5}
For the geodesic cutting and canonical capping operations of section 6, all boundary circle gluings are analytic diffeomorphisms. Their existence and uniqueness do not use a separate rectifiability conjecture. The same conclusion holds for smooth boundary diffeomorphisms, but mere differentiability of a homeomorphism without a nonvanishing derivative is not the same hypothesis.
\end{corollary}
\begin{proof}
Geodesic arclength coordinates are analytic, and the toric identification in those coordinates is $\theta\mapsto\theta_0-\theta$. It is an analytic diffeomorphism and quasisymmetric. Theorem \ref{A3-welding} therefore applies and its seams are removable. On a compact circle a smooth diffeomorphism and its inverse have bounded derivatives, so it is bi-Lipschitz and the same argument applies. A differentiable homeomorphism can have a derivative flat at a point, and no such bound follows from that weaker wording. This qualification separates the two readings without attributing extra regularity to the manuscript.
\end{proof}
